\subsection{Sprint 6 Deployment and CodeRAG Integration}

During Sprint 6, the deployment configuration was extended to support the final
service architecture with a separate CodeRAG service. The main infrastructure goal
was to make the project reproducible from a clean Docker state while keeping the
RAG implementation in its own repository.

CodeRAG was connected to the CodePilot repository as a Git submodule:

\begin{verbatim}
coderag
\end{verbatim}

The submodule points to:

\begin{verbatim}
https://github.com/GitFlameAI/GitFlame-CodeRAG.git
branch: sprint-6/http-rag-service
\end{verbatim}

This approach keeps the RAG source code outside the main repository history while
still allowing the complete demo stack to be cloned and started with the correct
CodeRAG version. A fresh clone should therefore use:

\begin{verbatim}
git clone --recurse-submodules https://github.com/GitFlameAI/Capstone-Project-GitFlame-CodePilot.git
\end{verbatim}

If the repository was already cloned, the submodule can be initialized with:

\begin{verbatim}
git submodule update --init --recursive
\end{verbatim}

The Docker Compose deployment was updated through:

\begin{verbatim}
backend/deploy/docker-compose.sprint2.override.yml
\end{verbatim}

The complete Sprint 6 stack includes:

\begin{verbatim}
frontend
backend
database
redis
ml-service
agent-engine
agent-worker
recommendation-service
coderag-database
rag-service
\end{verbatim}

The new \texttt{coderag-database} service uses the \texttt{pgvector/pgvector:pg16}
image and stores CodeRAG indexing data in a separate Docker volume. This avoids
mixing CodeRAG vector storage with the main CodePilot PostgreSQL database. The
new \texttt{rag-service} service is built from the \texttt{coderag} submodule and
exposes the CodeRAG HTTP API.

The Agent Engine receives the internal Docker Compose URL for CodeRAG through:

\begin{verbatim}
RAG_BASE_URL=http://rag-service:8004
\end{verbatim}

For local host-side checks, the service is available at:

\begin{verbatim}
http://localhost:8004
\end{verbatim}

The root \texttt{.env.example} file was updated with Sprint 6 RAG configuration:

\begin{verbatim}
RAG_PORT=8004
CODERAG_DATABASE_PORT=5433
RAG_BASE_URL=http://localhost:8004
RAG_API_KEY=
EMBEDDING_MODEL=jinaai/jina-embeddings-v2-base-code
RAG_USE_DENSE=true
RAG_CANDIDATE_TOP_K=50
RAG_MAX_CONTEXT_FILES=20
RAG_MAX_CHUNKS_PER_FILE=3
RAG_MAX_CONTEXT_TOKENS=12000
\end{verbatim}

The complete stack is started with:

\begin{verbatim}
docker compose \
  -f docker-compose.yml \
  -f backend/deploy/docker-compose.sprint2.override.yml \
  up -d --build
\end{verbatim}

A clean reproducibility check was performed by removing containers and volumes,
then starting the full stack again:

\begin{verbatim}
docker compose \
  -f docker-compose.yml \
  -f backend/deploy/docker-compose.sprint2.override.yml \
  down -v

docker compose \
  -f docker-compose.yml \
  -f backend/deploy/docker-compose.sprint2.override.yml \
  up -d --build
\end{verbatim}

After startup, Docker Compose reported all services as healthy:

\begin{verbatim}
agent-engine              healthy
agent-worker              healthy
backend                   healthy
coderag-database          healthy
database                  healthy
frontend                  healthy
ml-service                healthy
rag-service               healthy
recommendation-service    healthy
redis                     healthy
\end{verbatim}

The following health and readiness checks were verified:

\begin{verbatim}
GET http://localhost:8000/ready
{"components":{"agent_engine":"ready","redis":"ready","storage":"ready"},"service":"backend","status":"ready"}

GET http://localhost:8002/ready
{"status":"ready","model":"Qwen/Qwen3-Coder-30B-A3B-Instruct","version":"3.0.0"}

GET http://localhost:8004/health
{"status":"ok","database":"ready","version":"1.0.0"}

GET http://localhost:8001/health
{"status":"ok","service":"ml_service"}

GET http://localhost:8003/health
{"status":"ok","model":"Qwen/Qwen3-Coder-30B-A3B-Instruct"}

Redis check:
PONG
\end{verbatim}

The CodeRAG search endpoint was also checked with a valid request body. Since the
clean local database did not contain an indexed repository, the expected result was
an empty list:

\begin{verbatim}
POST http://localhost:8004/search
{"results":[]}
\end{verbatim}

The deployment notes, submodule instructions, environment variables, health checks,
RAG smoke-test scenario, and reproducibility commands are documented in:

\begin{verbatim}
infra/sprint6-rag-deployment-runbook.md
\end{verbatim}
